% Einheit 3 von 3 – Kreisteile (bank/kreis/e3.jsonl, zusammenbau v0.1)
%% TODO Zweigzeile Teil 1: was hier gelernt wird (Satz in Schülersprache aus dem Einheitstitel) – Einheit 3
\einheitenkopf[e3]{Einheit 3 von 3 · Kreisteile}
\zweigzeile{ab Kl. 7, je nach Buch bis Kl. 9 · P10}
\setcounter{aufgabe}{19}

%% TODO Titel: Ich-kann-Satz fehlt in der Bank (Platzhalter: Kettenname) – Nr. 20
%% TODO Anweisung: ein Satz über der ersten Teilaufgabe fehlt in der Bank – Nr. 20
\begin{aufgabe}{Kreisteile}
\begin{teile}
\teil Welcher Teil des Kreises ist grau? Kreuze an.\\ \kreuz{$\frac{1}{2}$}\\ \kreuz{$\frac{1}{3}$}\\ \kreuz{$\frac{1}{4}$}
\end{teile}
%% TODO Darstellung für schwach fehlt in der Bank (kreis-e3-k1-s1-v1; Abschnitt „Für schwache Schüler“)
\swz{Ein Kreisausschnitt hat den Mittelpunktswinkel $\alpha = 90^\circ$. Welcher Teil des Kreises ist das? Gib den Anteil als Bruch an.}{}
%% TODO Darstellung für schwach fehlt in der Bank (kreis-e3-k1-s1-v2; Abschnitt „Für schwache Schüler“)
\swz{Ein Kreisausschnitt hat den Mittelpunktswinkel $\alpha = 180^\circ$. Welcher Teil des Kreises ist das? Gib den Anteil als Bruch an.}{}
%% TODO Darstellung für schwach fehlt in der Bank (kreis-e3-k1-s1-v3; Abschnitt „Für schwache Schüler“)
\swz{Ein Kreisausschnitt hat den Mittelpunktswinkel $\alpha = 60^\circ$. Welcher Teil des Kreises ist das? Gib den Anteil als Bruch an.}{}
%% TODO Darstellung für schwach fehlt in der Bank (kreis-e3-k1-s1-v4; Abschnitt „Für schwache Schüler“)
\swz{Ein Stück einer runden Torte hat an der Spitze den Winkel $60^\circ$. Welcher Teil der Torte ist das? Gib den Anteil als Bruch an.}{}
%% TODO Darstellung für schwach fehlt in der Bank (kreis-e3-k1-s1-v5; Abschnitt „Für schwache Schüler“)
\swz{In einem Kreisdiagramm hat ein Sektor den Winkel $90^\circ$. Welcher Teil des Kreises ist das? Gib den Anteil als Bruch an.}{}
\swfrage{Was bleibt gleich, was ändert sich?}
%% TODO Darstellung für schwach fehlt in der Bank (kreis-e3-k1-s2-v1; Abschnitt „Für schwache Schüler“)
\swz{Ein Kreisausschnitt hat den Mittelpunktswinkel $\alpha = 100^\circ$. Wie viel Prozent des Kreises sind das? Runde auf eine Stelle.}{}
%% TODO Darstellung für schwach fehlt in der Bank (kreis-e3-k1-s3-v1; Abschnitt „Für schwache Schüler“)
\swz{Ein Kreisausschnitt ist $\frac{3}{8}$ des ganzen Kreises. Wie groß ist sein Mittelpunktswinkel?}{}
%% TODO Darstellung für schwach fehlt in der Bank (kreis-e3-k1-s4-v1; Abschnitt „Für schwache Schüler“)
\swz{Ein Kreisausschnitt hat den Radius $r = 9$ cm und den Mittelpunktswinkel $\alpha = 40^\circ$. Wie lang ist sein Bogen $b$? Rechne mit der $\pi$-Taste und runde auf eine Stelle.}{}
%% TODO Darstellung für schwach fehlt in der Bank (kreis-e3-k1-s5-v1; Abschnitt „Für schwache Schüler“)
\swz{Ein Kreisausschnitt hat den Radius $r = 6$ cm und den Mittelpunktswinkel $\alpha = 84^\circ$. Wie groß ist seine Fläche? Rechne mit der $\pi$-Taste und runde auf eine Stelle.}{}
%% TODO Darstellung für schwach fehlt in der Bank (kreis-e3-k1-s6-v1; Abschnitt „Für schwache Schüler“)
\swz{Ein Kreisausschnitt hat den Radius $r = 8$ cm und den Mittelpunktswinkel $\alpha = 45^\circ$. Wie groß ist der Umfang des Ausschnitts? Rechne mit der $\pi$-Taste und runde auf eine Stelle.}{}
%% TODO Darstellung für schwach fehlt in der Bank (kreis-e3-k1-s7-v1; Abschnitt „Für schwache Schüler“)
\swz{Ein Kreisring hat außen den Radius $7$ cm und innen den Radius $4$ cm. Wie groß ist seine Fläche? Rechne mit der $\pi$-Taste und runde auf eine Stelle.}{}
\swz[3]{Im Bild ist ein Kreisausschnitt grau gefärbt. Sein Mittelpunktswinkel beträgt $130^\circ$, der Winkel des weißen Teils heißt $\alpha$. Wie viel Prozent der Kreisfläche sind grau? Runde auf eine Stelle \hfill (P10 2025 OS).}{\kreissektor{130}{$130^\circ$}}
\end{aufgabe}

%% TODO Titel: Ich-kann-Satz fehlt in der Bank (Platzhalter: Kettenname) – Nr. 21
%% TODO Anweisung: ein Satz über der ersten Teilaufgabe fehlt in der Bank – Nr. 21
\begin{aufgabe}{Sachaufgabe (Rasensprenger, Tortenstück, Sektor im Kreisdiagramm)}
%% TODO Darstellung für schwach fehlt in der Bank (kreis-e3-k2-s1-v1; Abschnitt „Für schwache Schüler“)
\swz{Ein Rasensprenger spritzt $6$ m weit und dreht sich dabei hin und her über einen Winkel von $150^\circ$. Wie groß ist die Fläche, die er bewässert? Rechne mit der $\pi$-Taste und runde auf eine Stelle.}{}
\end{aufgabe}

%% TODO Titel: Ich-kann-Satz fehlt in der Bank (Platzhalter: Kettenname) – Nr. 22
%% TODO Anweisung: ein Satz über der ersten Teilaufgabe fehlt in der Bank – Nr. 22
\begin{aufgabe}{Kreisteile – Fehler finden, Begründen, Anwendung, Darstellungswechsel}
\swz[3]{Im Bild ist ein Ausschnitt mit $110^\circ$ grau, der Rest ist weiß. Jana berechnet den grauen Anteil: \rechnung{250^\circ : \text{Vollwinkel} &\approx 0{,}694 \\ &= 69{,}4\,\%} Wo ist der Fehler?}{\kreissektor{110}{$110^\circ$}}
\swfrage{Begründe, warum ein Halbkreis zum Mittelpunktswinkel $180^\circ$ gehört.}
\swz[3]{Ein Rasensprenger steht in der Ecke eines Gartens (rechter Winkel) und spritzt $7$ m weit. Auf jeden Quadratmeter sollen $4$ Liter Wasser fallen. Wie viel Wasser braucht er für einen Durchgang?}{}
\swz[3]{Das Kreisdiagramm zeigt, wie Schüler einer Klasse zur Schule kommen. Gib zu jedem Sektor den Anteil in Prozent an.}{\kreisdiagramm{Bus $162^\circ$/162, Rad $108^\circ$/108, Auto $90^\circ$/90}}
\end{aufgabe}

\uebersichtskasten{Kreisausschnitt: Der Mittelpunktswinkel sagt, welcher Anteil vom ganzen Kreis gemeint ist – Winkel geteilt durch 360°. \\ α = 72°: 72 : 360 = 1/5 = 20 \% \quad Bogen: 1/5 vom Umfang \quad Fläche: 1/5 von π · r² \\ r = 10 cm, α = 72°: b = 0,2 · 2 · π · 10 ≈ 12,6 cm \quad A = 0,2 · π · 10² ≈ 62,8 cm² \\ Umfang des Ausschnitts: Bogen plus zwei Radien. \quad 12,6 + 2 · 10 = 32,6 cm}
